\section{Two markings on one cyclotomic group}
For an odd prime $p$ and $n\geq2$, let $A=\mathbb Z[\zeta_p]/(\zeta_p-1)^n$,
$P=A\rtimes\gen\beta$, with $v^\beta=\zeta_pv$ and $|\beta|=p$.
Let $\sigma$ act by conjugation on $A$ and inversion on $\beta$;
let $\tau$ negate $A$ and fix $\beta$.
The shared split cyclotomic construction of~\cite{UF} proves that
$R=P\rtimes\gen{\sigma,\tau}$ has order $4p^{n+1}$ and the
marking
\[
(s\tau\sigma,\ \beta\tau\sigma,\ \sigma),\qquad s=1\in A,
\]
of type $\{p,2p\}$. The complete construction and its proof are included in the common
paper; the first application is not an additional proof premise. Taking $n=2,3$ gives respectively the first
two families in source Remark 1.3 for every odd prime.

There is a second marking on the same group. In
$G=P\rtimes\gen\sigma$ put $a=\beta s$ and $b=s^{-1}\sigma$.
The norm identity gives $|a|=p$; and $|b|=2|s|=2p^e$, where
$n=(e-1)(p-1)+r$, $1\leq r\leq p-1$.
Their product is $\beta\sigma$, an involution. They generate $G$:
$b^2=s^{-2}$ recovers $s$, then $\beta$ and $\sigma$.
Their cyclic subgroups intersect trivially: a common element lies
in $P\cap\gen b=\gen s$, whereas every nonidentity power of $a$
has nonzero image in $P/A$.
The element $c=\tau\sigma s$ in $R$ is an involution and lies
outside $G$. It fixes $\sigma$ by conjugation, inverts $s$, and
sends $\beta$ to $s^{-1}\beta^{-1}s$.
Consequently it simultaneously inverts $a,b$.
The rotation-pair argument of~\cite{UF} gives a second regular marking of
$R$ of type $\{p,2p^e\}$.
For $p=3$, $n=3$, this supplies the fourth claimed type
$\{3,18\}$ at order $324=4\cdot3^4$.

\section{A mixed Heisenberg construction}
For every odd prime $p$, define $H$ by triples
\[
(i,j,k)\in\mathbb Z/p^2\mathbb Z\times\mathbb Z/p\mathbb Z
                    \times\mathbb Z/p\mathbb Z,
\]
with multiplication
\[
(i,j,k)(i',j',k')=(i+i',j+j',k+k'-i'j).
\]
The formula is associative because its final correction is bilinear;
it has an identity and inverses. Set $x=(1,0,0)$, $v=(0,1,0)$,
$c=(0,0,1)$. Then $[x,v]=c$, $c$ is central, and
\[
|H|=p^4,\quad |x|=p^2,\quad |v|=|c|=p,
\quad H=\gen{x,v}.
\]
In particular this explicitly realizes the presentation
$x^{p^2}=v^p=c^p=[c,x]=[c,v]=1$, $[x,v]=c$;
its normal forms give the matching upper bound.
Define
\[
\theta(i,j,k)=(-i,j-i,-i(i+1)/2-k),
\]
where the last two coordinates are taken modulo $p$.
The value is independent of the representative of $i$, and direct
substitution in the bilinear multiplication proves it is an
automorphism. Alternatively the exact presentation shows it sends
$x\mapsto v^{-1}x^{-1}$, $v\mapsto v$, $c\mapsto c^{-1}$.
Its square is $(i,j,k)\mapsto(i,j,k+i)$, conjugation by $v$.

In the semidirect product of $H$ by an infinite cyclic generator $b$
acting by $h^b=\theta(h)$, the element $b^2v^{-1}$ is central:
$\theta^2$ is conjugation by $v$ and $\theta(v)=v$.
Quotient by its cyclic subgroup. The intersection of that subgroup
with $H$ is trivial, as its nontrivial powers have nonzero integer
coordinate. The quotient $G$ therefore contains $H$ with index two
and satisfies $b^2=v$; hence $|G|=2p^4$ and $|b|=2p$.
The relation $x^b=v^{-1}x^{-1}$ gives $(xb)^2=1$.
Moreover $G=\gen{x,b}$ and $\gen x\cap\gen b=1$, since the
intersection is inside $\gen x\cap\gen v=1$ in the explicit coordinates.
The presentation is invariant under $x\mapsto x^{-1}$,
$b\mapsto b^{-1}$: it inverts $v=b^2$, preserves the central
commutator of order $p$ and all displayed orders.
The resulting map is involutory. The common rotation-pair argument constructs a
regular polyhedron of order $4p^4$ and type $\{p^2,2p\}$.
This supplies the third family without a group catalogue.

\section{Why the fourth family has only the prime three}
For $p\geq5$, suppose a rank-three string C-group $R$ had order
$4p^4$ and type $\{p,2p^2\}$. Its Sylow $p$-subgroup $P$
is normal: its number divides four and is one modulo $p$.
The commuting endpoints generate a Klein four complement, so
$R/P\cong C_2^2$. Put $a=r_0r_1$ and $b=r_1r_2$.
Since $|a|=p$ is odd, $r_0P=r_1P$, whereas $r_2P$ is the other
independent endpoint image. The rotation subgroup $G=\gen{a,b}$
is normal of index at most two, because $r_1$ inverts both its
generators and $R=\gen{G,r_1}$. Its image in $R/P$ has order two,
so its index is exactly two. Thus $|G|=2p^4$.

As in the structural proof, $P=\gen{x,y}$ for $x=a$ and $y=b^2$.
It is nonabelian: an abelian group generated by these disjoint cyclic
subgroups would have order only $p^3$. The elementary small-kernel
lemma in the common paper says that
$P'$ has exponent $p$ and $P$ has class at most three.
For $p\geq5$, collection gives $(uv)^p=u^pv^p$ in $P$.
But $x$ and $y$ have orders $p,p^2$, and
$x^b=y^{-1}x^{-1}$. Taking $p$th powers contradicts $|y|=p^2$.
Thus the fourth family of Remark 1.3 does not exist for $p\geq5$;
the preceding cyclotomic example proves existence at $p=3$.
This corrects an auxiliary printed assertion, not a numbered open
problem, and makes no priority claim.